\chapter{Fine Needle Aspiration (FNA)}

\section*{Overview}
Fine needle aspiration (FNA) is a minimally invasive biopsy procedure in which a very thin needle is inserted into a lump or nodule to collect cells for microscopic examination, usually to determine whether a mass is benign or malignant.

\section*{Alternative Names}
FNA, fine-needle aspiration biopsy, FNA biopsy, aspiration cytology


\section*{Principle}
FNA uses a thin needle (22-27 gauge) to aspirate cells from a mass. Negative pressure is applied with a syringe. The aspirated material is expelled onto slides, smeared, and fixed (air-dried for Diff-Quik, ethanol-fixed for Papanicolaou). A portion may be placed in liquid medium for cell block preparation.

\section*{History}
Fine needle aspiration was developed in the early 20th century at Memorial Sloan Kettering. The technique was refined in Scandinavia in the 1960s and 1970s, becoming a standard minimally invasive diagnostic method.

\section*{Why It Is Done}
Your doctor orders an FNA when you have a lump or nodule that can be felt or seen on imaging, and they need to know whether it is cancerous. FNA is commonly used for thyroid nodules, lymph nodes, breast lumps, salivary gland lumps, and lung nodules seen on CT scan. It answers the question: does this lump contain cancer cells?

\section*{Is This a Primary or Secondary Test or Procedure}
Primary for many lumps and nodules. FNA is less invasive than a surgical biopsy and provides rapid answers.

\section*{How It Is Performed}
A very thin needle (25 to 27 gauge --- thinner than the needles used for blood draws) is inserted into the lump, often guided by ultrasound. The needle is moved back and forth within the lump to collect cells. The cells are then placed on a glass slide or in a liquid preservative and sent to the laboratory for examination. The test takes about fifteen minutes. Local anesthetic may be used, but many patients do not need it.

\section*{Preparation}
No special preparation is needed for most FNA procedures. If you take blood thinners, inform your doctor, though most FNAs can be performed safely.

\section*{Contraindications and Precautions}
This procedure has the following contraindications and precautions:

\begin{itemize}
    \item \textbf{Absolute contraindications}: Very few. Do not aspirate lesions suspected to be a pheochromocytoma or carotid body tumor without special precautions.
    \item \textbf{Relative contraindications}: Uncorrectable bleeding diathesis, severe thrombocytopenia, or therapeutic anticoagulation (hold anticoagulants when clinically safe); uncooperative patient; lesions deep or adjacent to major vascular structures without image guidance.
    \item \textbf{Precautions}: Use ultrasound or CT guidance for deep or small lesions; obtain informed consent about the small risk of a non-diagnostic sample; apply firm pressure after the procedure to minimize hematoma.
    \item \textbf{Special populations}: In pregnancy, avoid ionizing radiation by using ultrasound guidance; children may require general anesthesia or sedation for cooperation.
\end{itemize}

\section*{Risks}
This procedure carries the following risks:

\begin{itemize}
    \item \textbf{Bleeding or bruising}: Minor bleeding or bruising at the puncture site is common and usually self-limited.
    \item \textbf{Infection}: Very rare (<0.1\%) at the puncture site. Proper skin antisepsis minimizes this risk.
    \item \textbf{Pain or discomfort}: Brief pain during needle insertion. Some patients experience a vasovagal reaction (fainting, dizziness, nausea).
    \item \textbf{Hematoma}: Collection of blood under the skin if pressure is not held adequately after the procedure.
    \item \textbf{Pneumothorax}: Occurs when the needle enters the chest cavity (e.g., lung or hilar lymph node FNA); may require observation or chest tube drainage.
    \item \textbf{Tumor seeding}: Extremely rare (<0.01\%) spread of malignant cells along the needle tract.
    \item \textbf{Non-diagnostic sample}: Insufficient cellular material occurs in a minority of cases and may require a repeat FNA or a surgical biopsy.
\end{itemize}

\section*{Aftercare and Recovery}
You may have mild bruising at the site, which resolves within a few days. Results are typically available within one to three days.

\section*{Results and Interpretation}
FNA results are reported using established classification systems (e.g., the Bethesda System for thyroid nodules). A benign or negative result indicates no malignant cells, while atypical, suspicious, or malignant categories warrant further evaluation. Because FNA samples only a small amount of tissue, a negative result does not completely exclude cancer, and false negatives occur in a small percentage of cases.

\section*{Expected Results}
A normal (negative) result shows no malignant cells. May be reported as benign, atypical, suspicious, or malignant using the appropriate classification system.

\begin{center}
\begin{tabular}{ll}
\toprule
\textbf{Population} & \textbf{Reference Range} \\
\midrule
Adult Male & No malignant cells \\
Adult Female & No malignant cells \\
Children & Same as adult \\
Elderly & Same as adult \\
Pregnancy & Same as adult \\
\bottomrule
\end{tabular}
\end{center}
\section*{Follow-up}
Core needle biopsy, open surgical biopsy, or ultrasound surveillance.

\section*{References}
\begin{enumerate}
  \item MedlinePlus. Fine Needle Aspiration (FNA). U.S. National Library of Medicine; 2025.
  \item Current Medical Diagnosis and Treatment. McGraw-Hill; 2025.
\end{enumerate}

\clearpage
